% Общая преамбула конспекта: её подключают и общий файл Konspekt.tex, и файлы отдельных глав ch<N>/ch<N>.tex.
\usepackage[T2A]{fontenc}
\usepackage[utf8]{inputenc}
\usepackage[russian]{babel}
\usepackage{amsmath,amsthm,amssymb,mathtools}
\usepackage[margin=2.3cm]{geometry}
\usepackage{enumitem}
\usepackage{booktabs}
\usepackage{import}
\usepackage{tikz-cd}
\usetikzlibrary{babel}
\usetikzlibrary{decorations.pathreplacing,calc,arrows}
\colorlet{cnew}{red!75!black}   % стрелки и объекты, которые строятся в доказательстве
\colorlet{cold}{blue!65!black}  % данные
\newcommand{\figcap}[1]{\par\smallskip{\small\itshape #1}\par\medskip}
\usepackage{hyperref}
\hypersetup{colorlinks=true,linkcolor={blue!50!black},urlcolor={blue!50!black},linktoc=page,
  pdfauthor={},pdftitle={Поля и теория Галуа: конспект}}

% Главы нумеруются, разделы внутри глав — нет (номера вида 7.1 заняты утверждениями).
\setcounter{secnumdepth}{0}
\renewcommand{\theequation}{\arabic{equation}}

% Утверждения нумеруются вручную номерами книги: \begin{proposition}[7.2] или [7.13: основная теорема].
% Номер (часть до двоеточия) автоматически становится меткой st:<номер>; ссылка — \rf{<номер>}:
% если такое утверждение есть в собранном файле, номер становится переходом, иначе печатается как текст.
\makeatletter
\def\st@num#1:#2\@nil{#1}
\newcommand{\stmtanchor}[1]{\edef\st@tmp{\st@num#1:\@nil}%
  \phantomsection\edef\@currentlabel{\st@tmp}\label{st:\st@tmp}}
\DeclareRobustCommand{\rf}[1]{\@ifundefined{r@st:#1}{#1}{\hyperref[st:#1]{#1}}}
\makeatother

\newtheoremstyle{mplain}{\medskipamount}{\medskipamount}{\itshape}{}{\bfseries}{.}{ }{\thmname{#1}\thmnote{ #3\stmtanchor{#3}}}
\newtheoremstyle{mdef}{\medskipamount}{\medskipamount}{\normalfont}{}{\bfseries}{.}{ }{\thmname{#1}\thmnote{ #3\stmtanchor{#3}}}
\theoremstyle{mplain}
\newtheorem*{theorem}{Теорема}
\newtheorem*{proposition}{Предложение}
\newtheorem*{lemma}{Лемма}
\newtheorem*{corollary}{Следствие}
\theoremstyle{mdef}
\newtheorem*{definition}{Определение}
\newtheorem*{example}{Пример}
\newtheorem*{remark}{Замечание}
\newtheorem*{aside}{Отступление}
\newtheorem*{exercise}{Упражнение}

\DeclareMathOperator{\Gal}{Gal}
\DeclareMathOperator{\Aut}{Aut}
\DeclareMathOperator{\Hom}{Hom}
\DeclareMathOperator{\Homc}{Hom_{cont}}
\DeclareMathOperator{\Sym}{Sym}
\DeclareMathOperator{\Stab}{Stab}
\DeclareMathOperator{\Orb}{Orb}
\DeclareMathOperator{\Frob}{Frob}
\DeclareMathOperator{\sgn}{sgn}
\DeclareMathOperator{\id}{id}
\newcommand{\Q}{\mathbb Q}
\newcommand{\Z}{\mathbb Z}
\newcommand{\R}{\mathbb R}
\newcommand{\CC}{\mathbb C}
\newcommand{\F}{\mathbb F}
\newcommand{\Zh}{\widehat{\mathbb Z}}
\newcommand{\Zp}{\mathbb Z_p}
\newcommand{\Qp}{\mathbb Q_p}
\newcommand{\Fpb}{\overline{\mathbb F}_p}
\newcommand{\Qal}{\mathbb Q^{\mathrm{al}}}
\newcommand{\Qab}{\mathbb Q^{\mathrm{ab}}}
\newcommand{\Fsep}{F^{\mathrm{sep}}}
\newcommand{\varlim}{\varprojlim}
\newcommand{\Om}{\Omega}
